\documentclass[french,a4paper]{article}

\usepackage[utf8]{inputenc}
\usepackage[T1]{fontenc}
\usepackage{lmodern}
\usepackage{geometry}
\usepackage{graphicx}
\usepackage{babel}
\usepackage{amsmath}
\usepackage{amsthm}
\usepackage{amssymb}
\usepackage{hyperref}
\usepackage{stmaryrd}
\usepackage{enumitem}
\usepackage{tcolorbox}
\usepackage{dsfont}
\usepackage{multirow}
\usepackage{etoc}
\usepackage{titlesec}
\usepackage{esint}
\usepackage{cancel}
\usepackage{mathdots}

\setlist[itemize]{label=$\bullet$}


\renewcommand{\P}{\mathbb{P}}
\newcommand{\N}{\mathbb{N}}
\newcommand{\Z}{\mathbb{Z}}
\newcommand{\Q}{\mathbb{Q}}
\newcommand{\R}{\mathbb{R}}
\newcommand{\C}{\mathbb{C}}
\newcommand{\K}{\mathbb{K}}
\renewcommand{\L}{\mathbb{L}}
\newcommand{\U}{\mathbb{U}}
\newcommand{\st}{^*}
\newcommand{\tq}{\middle|}
\newcommand{\inv}{^{-1}}
\newcommand{\E}{\mathbb{E}}
\newcommand{\F}{\mathbb{F}}
\newcommand{\V}{\mathbb{V}}
\renewcommand{\ker}{\text{Ker}}
\newcommand{\im}{\text{Im}}
\newcommand{\card}{\text{Card}}
\newcommand{\Tr}{\text{Tr}}

\theoremstyle{plain}
\newtheorem{thm}{Théorème}
\newtheorem*{ind}{Indication}

\theoremstyle{definition}
\newtheorem{dfn}{Definition}

\theoremstyle{remark}
\newtheorem*{rmq}{Remarque}
\newtheorem*{app}{Application}

\newtcolorbox[auto counter]{ex}[2][]{colback=white, colframe=green!50!white, coltitle=black, fonttitle=\bfseries, title={Exercice~\thetcbcounter~: #2}, after title={\hfill #1}}

\title{TD Équations différentielles}

\begin{document}

\maketitle

\section{Résolution d'équations différentielles scalaires du premier ordre}
\section{Résolution d'équation différentielles scalaires d'ordre $2$}
\section{Systèmes différentiels}
\begin{ex}{$\Delta$ Sous-groupes à un paramètre de $\mathcal{GL}_n(\C)$}
	Soit $\varphi:\R\to\mathcal{GL}_n(\C)$ une morphisme de groupes continu.
	\begin{enumerate}
		\item Montrer que $\varphi$ est dérivable. On pourra introduire une primitive de $\varphi$.
		\item Montrer que $\varphi'=\varphi'(0)\varphi$.
		\item En déduire $\forall t\in\R,\ \varphi(t)=\exp(t\varphi'(0))$.
		\item $\star$ Décrire tous les morphismes de groupes continus de $\R$ dans $\mathcal{GL}_n(\C)$ puis de $\U$ dans $\mathcal{GL}_n(\C)$.
	\end{enumerate}
\end{ex}
\begin{proof}
	La quatrième question est difficile.
	\begin{enumerate}
		\item Comme d'habitude : intégrer par rapport à une variable en considérant $\Phi$ la primitive qui s'annule en $0$. Remarquer que $\Phi(u)=uI_n+o(u)$ donc puisque $\mathcal{GL}_n(\C)$ est ouvert, il existe $u_0$ non nul tel que $\Phi(u_0)$ soit inversible puis on conclut comme d'habitude.
		\item Dériver l'équation fonctionnelle par rapport à une des variables puis prendre la prendre égale à $0$.
		\item Invoquer une unicité à un problème de Cauchy.
		\item Pour la première partie de la question, c'est simplement la synthèse de l'analyse effectuée dans les trois question précédentes. Ensuite, on considère $f:\U\to\mathcal{GL}_n(\C)$ un morphisme de groupes continu. Alors $\varphi:t\mapsto f(e^{it})$ est un morphisme de groupes continu par composition. On raisonne ensuite par analyse-synthèse :
		\begin{itemize}
			\item Analyse : on peut fixer $A\in\mathcal{M}_n(\C)$ telle que $$\forall t\in\R,\ f(e^{it})=\varphi(t)=\exp(tA)$$ Or, $\exp(2\pi A)=f(1)=I_n$ donc $A$ est diagonalisable et $\text{Sp}(A)\subset i\Z$ (c'est un lemme qu'on démontre en utilisant la décomposition de Dunford, se référer au chapitre de réduction).
			\item Synthèse : réciproquement, on considère $$\begin{array}{lrcl}
				f:&\U&\longrightarrow&\mathcal{GL}_n(\C)\\
				&u&\longmapsto&\exp(tA)
			\end{array}$$ où $u=e^{it}$ et $A$ est diagonalisable à spectre inclus dans $i\Z$. On vérifie que cette fonction est bien continue en diagonalisant $A$ et en utilisant le fait que son spectre est inclus dans $i\Z$ et on vérifie immédiatement que c'est un morphisme de groupes continu.
		\end{itemize}
	\end{enumerate}
\end{proof}
\section{Étude d'équations différentielles du premier ordre}
\begin{ex}{$\Delta$ Équation différentielle vérifiée par le wronskien} On considère $n$ solutions de $y'=ay$ avec $a\in \mathcal{C}(I,\mathcal{L}(F))$ où $F$ est dimension finie, dont on prend $\mathcal{B}$ une base. On considère $y_1,\ldots, y_n$ des solutions de cette équation et $w$ leur wronskien. Soit $t_0\in I$. Alors : $$\forall t\in I,\ w(t)=w(t_0)\exp\left(\int_{t_0}^{t}\Tr(a(u))\ du\right)$$
\end{ex}
\begin{proof}
	En effet, d'après un exercice classique de formes $n$-linéaires alternées, on a pour tout $u$ et pour tout $t$ : $$\Tr(a(u))\det_\mathcal{B}(y_1(t),\ldots,y_n(t))=\sum_{j=1}^{n}\det_\mathcal{B}(y_1(t),\ldots,a(u)(y_j(t)),\ldots,y_n(t))$$ Soit, puisqu'on a des solutions de l'équation, $a(u)(y_j(t))=y_j'(t)$. On en déduit que $$\Tr(a(u))w(t)=\sum_{j=1}^{n}\det_\mathcal{B}(y_1(t),\ldots,y_j'(t),\ldots,y_n(t))=w'(t)$$ On en déduit la formule annoncée en résolvant l'équation différentielle.
\end{proof}
\begin{ex}{$\Delta$ $f'+f$ de limite nulle} Soif $f\in\mathcal{C}^1(\R_+,\R)$ telle que $f'+f\xrightarrow[+\infty]{}0$. Alors : $$\begin{cases}
		f\xrightarrow[+\infty]{}0 \\
		f'\xrightarrow[+\infty]{}0
	\end{cases}$$
\end{ex} \begin{proof}
	En effet, on pose $g=f'+f$ et on résout $y'+y=g$. Les fonctions $y$ solutions sont celles telles qu'il existe $C\in\R$ tel que $$\forall x\in\R_+,\ y(x)=e^{-x}\left(\int_{0}^{x}e^{t}g(t)\ dt +C\right)$$ Or, $f$ vérifie cette équation différentielle, donc il existe une telle constante $C$. De plus, puisque $e^{t}g(t)=o(e^{t})$, on a par intégration des relations de comparaion : $$\int_{0}^{x}e^{t}g(t)\ dt=o(e^x)$$ On en déduit que $f(x)=o(1)$, c'est-à-dire que $$f\xrightarrow[+\infty]{}0$$ Par différence avec $f'+f$, on obtient le résultat sur $f'$.
\end{proof}
\begin{app}
	Si $f$ est de classe $\mathcal{C}^2$ et vérifie $f''+3f'+2f\xrightarrow[+\infty]{}0$, alors $f$, $f'$ et $f''$ tendent vers $0$ quand $x$ tend vers $+\infty$. Il suffit d'appliquer deux fois ce qui précède à $g=f'+2f$.
\end{app}
\begin{rmq}
	Le résultat se généralise à $P(f)(x)\xrightarrow[x\to+\infty]{}0$ où $P$ est un polynôme complexe dont toutes les racines ont une partie réelle strictement négative. Le faire par récurrence en utilisant le cas de l'ordre $1$ déjà effectué.
\end{rmq}
\begin{ex}{$\star$ Moyenne de limite nulle}
	Soit $g:\R_+\to\R$ continue. On suppose que $$Mg:x\mapsto\frac{1}{x}\int_{0}^{x}g(t)\ dt$$ tend vers $0$ en $+\infty$. Montrer que pour toute solution $f$ de l'équation différentielle $y'+y=g$, la fonction $Mf$ tend vers $0$ en $+\infty$.
\end{ex}
\begin{proof}
	Exercice difficile. On cherche à faire comme dans l'exercice précédent. Soit $f$ une solution de cette équation. Alors : $$\forall x>0,\ \frac{f(x)-f(0)}{x}+Mf(x)=Mg(x)$$ Or, on remarque que $Mf$ est dérivable et que : $$\forall x>0,\ (Mf)'(x)=\frac{f(x)-f(0)}{x}-\frac{Mf(x)}{x}$$ Par conséquent, on a : $$(Mf)'(x)+\left(1+\frac{1}{x}\right)(Mf)(x)\xrightarrow[x\to+\infty]{}0$$ Ensuite, comme dans l'exercice précédent pour conclure, mais la résolution est légèrement plus difficile. On utilisera encore l'intégration des équivalents.
\end{proof}

\begin{ex}{$\Delta\star$ Lemme de Gronwall (HP)} Soit $I$ un intervalle de $\R$, $t_0\in I$, ainsi que que $a:I\to\R_+$ et $b:I\to\R_+$ deux fonctions continues. On note $A$ la primitive de $a$ s'annulant en $t_0$. Si $y$ de classe $\mathcal{C}^1$ vérifie : $$\forall t\geqslant t_0,\ \lVert y'(t)\rVert\leqslant b(t)+a(t)\lVert y(t)\rVert$$ alors $y$ vérifie : $$\forall t\geqslant t_0,\ \lVert y(t)\rVert\leqslant\lVert y(t_0)\rVert e^{A(t)}+e^{A(t)}\int_{t_0}^{t}b(s)e^{-A(s)}\ ds$$ Le résultat est facile à retrouver : on majore $\lVert y\rVert$ par la solution du problème de Cauchy : $$\begin{cases}
		z'=b(t)+a(t)z \\
		z(t_0)=\lVert y(t_0)\rVert
	\end{cases}$$
\end{ex}
\begin{proof}
	La concept est simple : on fait des choses sans comprendre pourquoi, mais ça marche et tout va se simplifier par miracle. \textit{cf} le cours pour les détails.
\end{proof}
\begin{app}
	Étude du comportement asymptotique des solutions de l'équation : $$y''+\frac{1}{x^2}y'+y=0$$
\end{app}
\begin{ex}{$\star$ Une application du lemme de Gronwall}
	Soit $A\in\mathcal{C}^0(\R_+,\mathcal{M}_n(\R))$ intégrable sur $\R_+$. On s'intéresse à l'équation différentielle $(E)$ : $$Y'=A(t)Y$$
	\begin{enumerate}
		\item Montrer que toutes les solutions de $(E)$ sont bornées puis qu'elles admettent une limite finie en $+\infty$.
		\item Montrer que l'application $u$ qui à un vecteur $Y_0$ associe la limite en $+\infty$ de la solution de $(E)$ valant $Y_0$ en $0$ est un automorphisme de $\R^n$.
	\end{enumerate}
\end{ex}
\begin{proof}
	Exercice difficile.
	\begin{enumerate}
		\item Appliquer le lemme de Gronwall pour montrer que $Y$ est bornée puisque $A$ est intégrable. Ainsi, $Y'=O(A)$ donc $Y'$ est intégrable donc $Y$ admet une limite finie en $+\infty$.
		\item Deux méthodes sont possibles :
		\begin{itemize}
			\item une méthode stylée mais utilisant d'autres exercices : on considère $w$ le wronskien des solutions qui valent les vecteurs de la base canonique en $0$. Un exercice assure que $$w(t)=\underbrace{w(0)}_{=1}\exp\left(\int_{0}^{t}\Tr(A(s))\ ds\right)$$ Or, cette intégrale converge car la trace est linéaire et $A$ est intégrable. Donc le wronskien admet une limite finie non nulle (car il y a une exponentielle). Or, par continuité du déterminant, cette limite est le déterminant de $u$. Donc $u$ est inversible.
			\item une méthode faisable sans avoir besoin d'un autre exo : on constate que $$u(Y_0)-Y(t)=\int_{t}^{+\infty}A(s)Y(s)\ ds$$ Ensuite, on montre que $u$ est injective ce qui suffira à conclure par égalité de dimensions. Supposons $u(Y_0)=0$. Alors : $$\lVert Y(t)\rVert\leqslant \int_{t}^{+\infty}\lVert A(s)\rVert\lVert Y(s)\rVert\ ds$$ Prenons $t_0$ tel que $$\int_{t_0}^{+\infty}\lVert A(s)\rVert\ ds <1$$ (possible car c'est le reste d'une intégrale convergente). Alors : $$\lVert Y\rVert_{\infty,[t_0,+\infty[}\leqslant\lVert Y\rVert_{\infty,[t_0,+\infty[}\int_{t_0}^{+\infty}\lVert A(s)\rVert\ ds$$ donc $$\lVert Y\rVert_{\infty,[t_0,+\infty[}=0$$ On en déduit aisément que $Y=0$, car par exemple $Y'(t_0+1)=0$ et $Y(t_0+1)=0$. Donc $u$ est injective ce qui conclut.
		\end{itemize}
	\end{enumerate}
\end{proof}
\begin{ex}{$\star$ Une équation différentielle perturbée}
	Soit $A\in\mathcal{M}_n(\C)$ diagonalisable à valeur propres imaginaires pures et $B:\R_+\to\mathcal{M}_n(\C)$ continue et intégrable. On considère une solution non nulle $X$ de l'équation différentielle $$X'=(A+B(t))X$$ Montrer qu'il existe un vecteur non nul $X_0$ tel que $$X(t)=\exp(tA)X_0+o(1)$$ au voisinage de $+\infty$.
\end{ex}
\begin{proof}
	Exercice difficile. D'après le lemme de Gronwall, $X$ est bornée. En écrivant $X$ sous la forme $X=e^{tA}Y$ on trouve $$\forall t\in\R_+,\ X(t)=e^{tA}\underbrace{\left(Y_0+\int_{0}^{t}\underbrace{e^{-sA}B(s)X(s)}_{=O(B(s))}\ ds\right)}_{\xrightarrow[t\to+\infty]{}X_0}$$ donc $X(t)=\exp(tA)X_0+o(1)$ au voisinage de $+\infty$. Il reste à montrer que $X_0\ne 0$. Deux pistes (j'ai essayé la première) : 
	\begin{itemize}
		\item Remarquer que l'application qui à $X$ associe $X_0$ est linéaire et tenter d'exhiber sa bijection réciproque en perturbant l'équation différentielle à l'envers ;
		\item Supposer que $X_0=0$ et montrer que $X=0$.
	\end{itemize} A finir...
\end{proof}
\begin{ex}{$\Delta$ Théorème de Floquet}
	Soit $A\in\mathcal{C}^0(\R,\mathcal{M}_n(\C))$ périodique de période $T$. montrer que l'équation différentielle $$Y'=A(t)Y$$ admet une solution non nulle $f$ pour laquelle il existe $\lambda\in\C$ tel que $$\forall t\in\R,\ f(t+T)=\lambda f(t)$$
\end{ex}
\begin{proof}
	Vérifier que l'application qui à $Y$ une solution de cette équation associe $t\mapsto Y(t+T)$ est un endomorphisme de $S$ l'espace des solutions. il suffit alors de considérer un vecteur propre et une valeur propre (possible car on travaille dans $\C$ donc cela existe).
\end{proof}
\begin{ex}{$\Delta$ Pas de sous-espace de dimension finie impaire stable}
	On considère l'application $T$ définie sur $\mathcal{C}^{\infty}(\R,\R)$ par : $$\forall x\in\R,\ T(f)(x)=f(x)+\int_{0}^{x}f(t)\ dt$$ Montrer que $f$ est un automorphisme de $\mathcal{C}^\infty(\R,\R)$. Existe-t-il un sous-espace vectoriel de $\mathcal{C}^\infty(\R,\R)$ de dimension finie impaire stable par $T$ ?
\end{ex}
\begin{proof}
	La caractère linéaire est immédiat. Pour l'injectivité et la surjectivité, utiliser le théorème de Cauchy linéaire. Pour la dernière question, on raisonne par l'absurde en supposant que c'est la cas. Alors, l'induit de $T$ sur cet espace admet une valeur propre et un vecteur propre associé. On distingue deux cas : si la valeur propre vaut $1$ ou non. Dans tous les cas, on montre que le vecteur propre est nul, ce qui est absurde. On en déduit qu'il n'existe pas de sous-espace de dimension finie impaire stable par $T$.
\end{proof}
\begin{ex}{Paramètre symétrique constant}
	Soit $A\in\mathcal{S}_n(\R)$. On considère l'équation différentielle $X'=AX$.
	\begin{enumerate}
		\item Soit $X$ une solution non nulle. Montrer que $t\mapsto\ln\left(\lVert X(t)\rVert_2\right)$ est convexe sur $\R$.
		\item On suppose de plus $A\in\mathcal{S}_n^{++}(\R)$. Soit $X$ une solution non nulle. Montrer que $t\mapsto \lVert X(t)\rVert_2$ établit une bijection de $\R$ sur $\R_+\st$.
	\end{enumerate}
\end{ex}
\begin{proof}
	On réalise tout d'abord un travail préliminaire utile pour les deux question. En effet, en notant $\lambda_i$ les éléments du spectre de $A$, on peut, après orthodiagonalisation, écrire $$\forall t\in\R, X(t)=P\text{Diag}(e^{\lambda_i t})P\inv X_0$$ Si $X$ est non nulle, $X_0$ est non nul et on pose $Y_0=P\inv X_0$ qui est donc tout aussi non nul. De plus, puisque $P$ est orthogonale, on sait que $$\forall t\in\R,\ \lVert X(t)\rVert =\lVert \text{Diag}(e^{\lambda_i t})Y_0\rVert$$ Par conséquent, après produit et expression de la norme $2$, on trouve : $$\forall t\in\R,\ \lVert X(t)\rVert=\sqrt{\sum_{i=1}^{n}e^{2\lambda_i t}v_i^2}$$
	Puis :
	\begin{enumerate}
		\item Passer au logarithme et montrer que la fonction obtenue (sans le 1/2) est convexe. Pour cela, la dériver deux fois et utiliser l'inégalité de Cauchy-Schwarz pour montrer que le numérateur est positif.
		\item Remarquer que les valeurs propres sont alors strictement positives. La fonction $t\mapsto\lVert X(t)\rVert$ est alors strictement croissante et ses limites sont $0$ en $-\infty$ et $+\infty$ en $+\infty$ puisque la solution est non nulle. Le théorème des valeurs intermédiaires et la stricte croissance permettent de conclure.
	\end{enumerate}
\end{proof}

\begin{ex}{Paramètre à valeurs antisymétriques}
	On considère l'équation $X'=A(t)X$ d'inconnue matricielle avec $\forall t\in\R,\ A(t)\in\mathcal{A}_n(\R)$. Soit $X$ une solution de cette équation différentielle. Montrer que si $X(0)\in\mathcal{O}_n(\R)$, alors $$\forall t\in\R,\ X(t)\in\mathcal{O}_n(\R)$$
\end{ex}
\begin{proof}
	Supposons $X(0)\in\mathcal{O}_n(\R)$ et considérons $f:t\mapsto X(t)^TX(t)$. Il est clair que $f$ est dérivable et : $$\forall t\in\R, f'(t)=X(t)^TA(t)X(t)+X(t)^TA(t)^TX(t)=X(t)^TA(t)X(t)-X(t)^TA(t)X(t)=0$$ si bien que $f$ est constante. Or, $f(0)=I_n$ puisque $X(0)\in\mathcal{O}_n(\R)$. Par conséquent : $$\forall t\in\R,\ f(t)=I_n$$ Autrement dit, on a bien : $$\forall t\in\R,\ X(t)\in\mathcal{O}_n(\R)$$
\end{proof}
\begin{ex}{$\star$ Une identité remarquable}
	On considère $A$ et $B$ deux matrices de $\mathcal{M}_n(\C)$ et on pose $C=AB-BA$. On suppose que $A$ et $B$ commutent avec $C$. On définit : $$\forall t\in\R,\ M(t)=e^{-t(A+B)}e^{tA}e^{tB}$$ Exprimer $M(t)$ en fonction de $C$ et $t$.
\end{ex}
\begin{proof}
	Exercice difficile. On utilise l'astuce de dérivation : en mettant la dérivée de la première fonction à sa droite, et celles des deux suivantes à leur gauche, on obtient : $$\forall t\in\R,\ M'(t)=e^{-t(A+B)}\left(-Be^{tA}+e^{tA}B\right)e^{tB}$$ On écrit alors : $$\forall t\in\R,\ -Be^{tA}+e^{tA}B=\sum_{n=0}^{+\infty}\frac{t^n}{n!}(A^nB-BA^n)$$ Or, il est classique, par récurrence, que $$\forall n\in\N\st,\ A^nB-BA^n=nAc^{n-1}$$ Par conséquent, après réindexation et utilisation de la définition de l'exponentielle, on en déduit : $$\forall t\in\R, M'(t)=e^{-t(A+B)}tCe^{tA}e^{tB}$$ Or, comme $C$ commute avec $A$, $C$ commute avec $e^{tA}$ (par continuité du produit), si bien qu'on a finalement : $$\forall t\in\R,\ M'(t)=tCM(t)$$ Comme de plus $M(0)=I_n$, on en déduit finalement l'identité : $$\forall t\in\R,\ M(t)=\exp\left(\frac{t^2}{2}C\right)$$ 
\end{proof}
\begin{ex}{Équation de Lax}
	Soit $B:\R\to\mathcal{M}_n(\C)$ une application continue et $A_0\in\mathcal{M}_n(\C)$. \begin{enumerate}
		\item Montrer qu'il existe $A:\R\to\mathcal{M}_n(\C)$ dérivable telle que $$\forall t\in\R,\ A'(t)=A(t)B(t)-B(t)A(t)$$ et $A(0)=A_0$.
		\item $\star$ Montrer qu'alors le spectre de $A(t)$ est indépendant de $t$, et plus précisément qu'il existe $S:\R\to\mathcal{GL}_n(\C)$ de classe $\mathcal{C}^1$ telle que $$\forall t\in\R,\ A(t)=S(t)A_0S(t)\inv$$
	\end{enumerate}
\end{ex}
\begin{proof}
	La première question est une application directe du cours, la deuxième est difficile.
	\begin{enumerate}
		\item Interpréter comme un problème de Cauchy avec une application linéaire.
		\item Chercher des infos sur $S$ par $CN$. Il suffit qu'elle vérifie $$S'=-BS$$ et $S(0)=I_n$. Réciproquement, on vérifie qu'une telle solution est de classe $C^1$ et convient, en démontrant la formule de la dérivée de l'inverse. Attention : il faut vérifier en appliquant le théorème de Cauchy linéaire que comme $S(0)\in\mathcal{GL}_n(\C)$, on a alors : $$\forall t\in\R,\ S(t)\in\mathcal{GL}_n(\C)$$
	\end{enumerate}
\end{proof}

\section{Étude d'équations différentielles d'ordre $2$}
\begin{ex}{$\Delta\star$ Méthode de Sturm-Liouville (HP)} Soit $p_1$ et $p_2$ des fonctions continues de $\R$ dans $\R$ telles que $p_1\leqslant p_2$. On considère les équations différentielles : $$\begin{cases}
		y''+p_1y=0 \ \ (E_1) \\
		y''+p_2y=0 \ \ (E_2)
	\end{cases}$$ Soit $y_1$ une solution réelle non nulle de $(E_1)$ et $y_2$ une solution de $(E_2)$. Si $t_1$ et $t_2$ sont deux zéros distincts de $y_1$, alors il existe $t\in[t_1,t_2]$ tel que $y_2(t)=0$.
\end{ex}
\begin{proof}
	Considérer $W=y_1y_2'-y_1'y_2=y_1y_2(p_1-p_2)$. Les zéros de $y_1$ sont isolés, donc on peut prendre SPG $t_1$ et $t_2$ deux zéros consécutifs de $y_1$. SPG, on peut supposer que $y_1$ est strictement positive entre $t_1$ et $t_2$ : on a alors $y_1'(t_1)>0$ et $y_2'(t_2)<0$ en effectuant un DL. On raisonne alors par l'absurde en supposant que $y_2$ ne s'annule pas sur $[t_1,t_2]$. Le Wronskien est alors monotone, mais s'il est croissant par exemple, alors il est positif en $t_1$ et négatif en $t_2$. Mais par conséquent, $y_1'(t_1)=0$. Donc $y_1$ est nulle car c'est une solution au problème de Cauchy : $$\begin{cases}
		y''+p_1y=0 \\
		y(t_1)=0\\
		y_1'(t_1)=0
	\end{cases}$$ C'est absurde, donc $y_2$ s'annule sur $[t_1,t_2]$. En examinant la preuve ci-dessus, on peut même montrer en réalité que $y_2$ s'annule dans $[t_1,t_2[$ et dans $]t_1,t_2]$.
\end{proof}
\begin{app}
	Application à l'étude des zéros d'une solution non nulle sur $\R_+$ de $y''+e^ty=0$. Montrer que ces zéros forment une suite strictement croissante, dont on déterminera un équivalent à l'infini.
\end{app}
\begin{app}
	Permet de montrer qu'une solution de $y''+qy=0$ avec $q>0$ et $q'>0$ au voisinage de l'infini admet une infinité de zéros.
\end{app}
\begin{ex}{$\star$ Une application de la méthode de Sturm-Liouville}
	Montrer que les zéros d'une solution $y$ sur $\R_+$ non nulle de $y''+e^ty=0$ forment une suite strictement croissante $(t_n)$. Déterminer un équivalent de la suite $(t_n)$.
\end{ex}
\begin{proof}
	Pour le fait qu'il y en a une infinité, utiliser la méthode de Sturm-Liouville avec $t\mapsto \sin(t)$ solution de $y''+y=0$ car $\forall t\in\R_+,\ e^t\geqslant 1$. Ensuite, pour les ordonner, on peut remarquer qu'il ne peut y avoir une infinité de zéros sur le compact $[0,A]$ grâce au principe des zéros isolés. On peut alors les ordonner. Pour l'équivalent, on utilise le fait que $$\forall t\in[t_n,t_{n+1}],\ e^{t_n}\leqslant e^t\leqslant e^{t_{n+1}}$$ On applique la méthode de Sturm-Liouville à $t\mapsto \sin(e^{t_n/2}(t-t_n))$, $y$ et $t\mapsto\sin(e^{t_{n+1}/2}(t-t_n))$ pour obtenir l'encadrement : $$\frac{\pi}{e^{t_{n+1/2}}}\leqslant t_{n+1}-t_n\leqslant\frac{\pi}{e^{t_n/2}}$$ Or, $t_n\to+\infty$ puisque par construction on a des zéros aussi grand que l'on veut (ou alors, par l'absurde, $t_n$ convergerait vers un zéro de $y$ par continuité, mais cela est impossible par stricte croissance de $(t_n)$). Par conséquent, $t_{n+1}-t_n\to 0$. On en déduit que $$t_{n+1}-t_n\sim \frac{\pi}{e^{t_n/2}}$$ car les exponentielles sont équivalentes car la différence des arguments tend vers $0$. On étudie ensuite la suite annexe $e^{t_n/2}$. En faisant la différence de deux termes, on trouve une limite non nulle avec les résultats précédents ($\pi/2$) donc on peut sommer les relations de comparaison. On en déduit que $$e^{t_n}\sim \frac{n^2\pi^2}{4}$$ et comme tout tend vers $+\infty$, on sait qu'on peut prouver rapidement qu'on peut passer les équivalents au $\ln$, ce que l'on fait et qui fournit finalement : $$t_n\sim 2\ln(n)$$
\end{proof}
\begin{ex}{Au moins une annulation}
	Soit $p\in\mathcal{C}(\R,\R_+\st)$ et $y$ une solution réelle sur $\R$ de l'équation $y''+p(x)y=0$. Montrer que $y$ s'annule au moins une fois.
\end{ex}
\begin{proof}
	Raisonnons par l'absurde. Alors, $y$ est de signe constant, SPG strictement positive. Donc $y$ est concave. Or, pour tout $x_0$, la dérivée de $y$ en $x_0$ ne peut pas être non nulle car alors $y$ s'annulerait par concavité comme $y$ est définie sur $\R$, donc $y'=0$. Donc $y$ est constante. Donc $y$ est nulle car $y''$ est nulle et $p$ ne s'annule jamais. C'est absurde.
\end{proof}
\begin{ex}{Une unique solution périodique}
	Soit $g$ une fonction continue périodique de $\R$ dans $\R$. Montrer que l'équation $$y''-3y'+2y=g(t)$$ admet exactement une solution périodique.
\end{ex}
\begin{proof}
	Pour l'unicité, il y a même unicité de la solution bornée en observant la forme des solutions de l'équation homogène, donc \textit{a fortiori} unicité d'une solution périodique. Pour l'existence, faire à moitié la variation des constantes (les intégrales n'importent pas) pour se ramener à un système $2\times 2$ \textit{via} la CNS habituelle pour les solutions périodiques (détaillée dans la section suivante).
\end{proof}
\begin{ex}{$\star$ Équation $y''+q(x)y$ avec $q$ positive croissante asymptotiquement} On considère $q\in\mathcal{C}^1(\R_+,\R)$ et l'équation $y''+q(x)y=0$. Alors :
	\begin{itemize}
		\item toute solution admet une infinité de zéros ;
		\item $\star\star$ toute solution est bornée
	\end{itemize}
\end{ex}
\begin{proof}
	Pour le premier point, appliquer la méthode de Sturm-Liouville sur un voisinage de $+\infty$ car on a alors $q(x)\geqslant k>0$ par croissance et positivité asymptotiques de $q$. On se ramène alors aux zéros de $\sin$ (à peu de choses près) que l'on contrôle bien, et la méthode de Sturm-Liouville prouver que toute solution de la première équation admet un zéro entre deux zéros de ce $\sin$. \\ Pour le deuxième point, considérer $z(x)=\displaystyle (y(x))^2+\frac{(y'(x))^2}{q(x)}$. Cette fonction est dérivable au voisinage de $+\infty$, et sa dérivée est négative (on a une simplification car $y$ est solution). Ainsi, $z$ est décroissante à partir d'un certain seuil, puis par positivité, $y^2$ est bornée dans ce voisinage de $+\infty$, donc $y$ y est bornée. Par continuité, $y$ est donc bornée sur $\R_+$ tout entier.
\end{proof}
\begin{ex}{$\Delta$ Existence de solutions non bornées} Considérons l'équation différentielle $y''+q(x)y=0$ avec $q$ continue et intégrable. Alors :
	\begin{itemize}
		\item La dérivée de toute solution bornée de cette équation tend vers $0$ en $+\infty$ ;
		\item En déduire qu'il existe des solutions à cette équation non bornées.
	\end{itemize}
\end{ex}
\begin{proof}
	Soit $y$ une solution de cette équation. Prouvons les deux points dans l'ordre.
	\begin{itemize}
		\item En intégrant $y''$ entre $0$ et $t$ et en remarquant que $qy$ est intégrable car c'est un $O(q)$, on trouve que $y'$ admet une limite finie en $+\infty$. Cette limite ne peut être non nulle, car en intégrant $y'$ au voisinage de $+\infty$ là où $y'$ serait de signe constant est de module minorée par un réel strictement positif, on trouverait que $y$ est non bornée.
		\item S'il existe une solution non nulle bornée $y_1$, on considère une deuxième solution $y_2$ telle que $(y_1,y_2)$ soit un système fondamental. Alors, le wronskien de $y_1$ et $y_2$ est constant car l'équation différentielle ne possède pas de terme en $y'$, et il est donc égal à une constante non nulle car $(y_1,y_2)$ est un système fondamental. De plus, si $y_2$ est bornée, alors d'après le premier point, le wronskien doit tendre vers $0$ en l'infini, donc être nul. Or, il est non nul, donc $y_2$ et non bornée. Dans tous les cas, il existe une solution non bornée.
	\end{itemize}
\end{proof}
\begin{ex}{$\star\star$ Inégalité de Liapounov}
	Soit $a<b$ et $f:[a,b]\to\R$. Soit $u$ une fonction non nulle de classe $\mathcal{C}^2$ solution de $y''+fy=0$ sur $[a,b]$ vérifiant $u(a)=u(b)=0$. Montrer l'inégalité de Liapounov : $$\int_{a}^{b}\lvert f\rvert \geqslant\frac{4}{b-a}$$
\end{ex}
\begin{proof}
	Puisque les zéros sont isolés, on peut sans perte de généralité supposer que $a$ et $b$ sont des zéros consécutifs de $u$, ce qui revient à prouver une inégalité encore plus forte que celle demandée. Ensuite, on utilise la méthode de variations des constantes pour montrer que : $$\forall t\in[a,b],\ (b-a)u(t)=(t-a)\int_{t}^{b}u(s)f(s)(b-s)\ ds+(b-t)\int_{a}^{t}u(s)f(s)(s-a)\ ds$$ Ensuite, on considère $t_0$ en lequel la norme infinie de $u$ sur $[a,b]$ est atteinte (qui est strictement positive car $u$ est non nulle). L'inégalité triangulaire permet d'obtenir : $$(b-a)\lVert u\rVert_\infty\leqslant (t_0-a)\int_{t_0}^{b}\lVert u\rVert_\infty \lvert f(s)\rvert (b-t_0)\ ds+(b-t_0)\int_{a}^{t_0}\lVert u\rVert_\infty \lvert f(s)\rvert (t_0-a)\ ds$$ On en déduit en regroupant les termes et en simplifiant par la norme infinie non nulle : $$b-1\leqslant (b-t_0)(t_0-a)\int_{a}^{b}\lvert f\rvert$$ Or, $x\mapsto (b-x)(x-a)$ atteint son maximum en $\displaystyle\frac{a+b}{2}$ (fonction polynomiale de degré 2), où elle vaut alors $\displaystyle\frac{(b-a)^2}{4}$. En majorant puis divisant par cette dernière quantité, on obtient bien : $$\int_{a}^{b}\lvert f\rvert \geqslant\frac{4}{b-a}$$
\end{proof}
\begin{ex}{$\star$ Théorème de stabilité de Liapounov}
	On considère $T>$ et $q$ une fonction continue sur $\R$ $T$-périodique. On note $S$ l'ensemble des solutions de l'équation $(E)$ : $$y''+qy=0$$ On note $y_1$ (resp. $y_2$) l'élément de $S$ vérifiant $y_1(0)=1$ et $y_1'(0)=0$ (resp. $y_2(0)=0$ et $y_2'(0)=1$).
	\begin{enumerate}
		\item Montrer que si $f\in S$, alors $f_T:x\mapsto f(x+T)$ appartient aussi à $S$. On note alors $\Phi$ l'endomorphisme de $S$ qui à $f$ associe $f_T$ ainsi que $A$ sa matrice dans la base $(y_1,y_2)$.
		\item Calculer le déterminant de $A$.
		\item On suppose que $\lvert\Tr(A)\rvert<2$. Montrer que tout élément de $S$ est borné.
		\item On suppose $q$ positive non identiquement nulle. Montrer que tout élément de $S$ possède possède au moins deux zéros.
		\item On suppose $q$ positive non identiquement nulle et $T\displaystyle\int_{0}^{T}q<4$. En utilisant l'inégalité de Liapounov, montrer que les solutions de $(E)$ sont toutes bornées.
	\end{enumerate}
\end{ex}
\begin{proof}
	Remarquons qu'on a pour tout $f\in S$, $f=f(0)y_1+f'(0)y_2$.
	\begin{enumerate}
		\item Immédiat en utilisant la $T$-périodicité de $q$.
		\item D'après notre remarque préliminaire, $\Phi(y_1)=y_1(T)y_1+y_1'(T)y_2$ et symétriquement pour $\Phi(y_2)$. Alors : $$A=\begin{pmatrix}
			y_1(T)&y_2(T)\\
			y_1'(T)&y_2'(T)
		\end{pmatrix}$$ Pour le déterminant de $A$, on reconnaît le Wronskien évalué en $T$. Or, le Wronskien est constant puisqu'il n'y a pas de terme en $y'$ dans $(E)$. Par conséquent, il est égal au Wronskien en $0$, soit : $$\det(A)=1$$
		\item Matrice $2\times 2$, on nous a demandé de calculer le déterminant à la question d'avant et on nous donne une information sur la trace, il faut penser polynôme caractéristique ! L'hypothèse nous assure que $\chi_A$ possède deux valeurs propres complexes non réelles conjuguées $\lambda$ et $\overline{\lambda}$. De plus, elle sont de module $1$ car leur produit vaut $\det(A)=1$. On considère $(u_1,u_2)$ une base de vecteurs propres des solutions complexes associée aux valeurs propres $(\lambda,\overline{\lambda})$. Alors, $\lvert u_1\rvert$ et $\lvert u_2\rvert$ sont $T$-périodiques. Elles sont donc bornées. Donc toute solution complexe est bornée. \textit{A fortiori}, toute solution réelle est bornée.
		\item En raisonnant un première fois par l'absurde en supposant qu'une solution $y$ ne s'annule pas, il est classique, en utilisant la convexité ou la concavité, que $y$ doit être constante, puis constante nulle. En raisonnant par l'absurde une seconde fois en supposant que $y$ possède uniquement un point d'annulation $x_0$, on peut sans perte de généralité supposer $y$ négative avant $x_0$ puis positive après. Alors, $y$ est concavec après $x_0$, donc $y'$ admet une limite dans $\R\cup\left\{-\infty\right\}$. Cette limite ne peut être strictement sinon positive sans quoi $y$ tendrait vers $-\infty$ en $+\infty$, ce qui contredirait sa positivité. Donc cette limite est positive et $y'$ est positive après $x_0$ donc $y$ croît après $x_0$. Elle admet donc une limite $l$ strictement positive en $+\infty$. Mais, à partir d'un certain $y''\leqslant\displaystyle-\frac12 lq$, et puisque $q$ n'est pas identiquement nulle et est périodique, l'intégrale de $y''$ diverge vers $-\infty$, donc $y'$ tend vers $-\infty$, ce qui est là encore absurde.
		\item On essaye de se ramener à la question $3$. En raisonnant par l'absurde, on suppose que $\lvert\Tr(A)\rvert\geqslant2$. Alors $A$ admet une valeur propre réelle $\lambda$, et on peut considérer $u$ un vecteur propre de l'endomorphisme $\Phi$ associé à la valeur propre $\lambda$. D'après la question 4, $u$ admet un zéro $a$. Puisque c'est un vecteur propre, $a+T$ est aussi un zéro de $u$. L'inégalité de Liapounov fournit alors : $$T\int_a^{a+T} q\geqslant 4$$ ce qui est absurde par hypothèse (la valeur de l'intégrale de $q$ sur une période est une constante par périodicité). On en déduit que $\lvert\Tr(A)\rvert<2$, puis la question 3 conclut.
	\end{enumerate}
\end{proof}
\begin{ex}{Ordre 2 et fausse périodicité}
	Soit $T>0$ ainsi que $(a,b)\in\R^2$. On considère $y$ une solution de l'équation différentielle $$y''+ay=b$$ telle que $y(T)=y(0)=0$. Montrer que $\forall n\in\Z,\ y(nT)=0$.
\end{ex}
\begin{proof}
	Considérer $z:t\mapsto y(t+T)$ la translation. Elle est aussi solution de l'équation. On considère alors le wronskien de $y$ et $z$. Celui-ci est constant car il n'y a pas de terme en $y'$ dans l'équation. Et, en évaluant en 0, il vaut $0$ donc est constamment nul. Si $y$ est la fonction nulle, le résultat est évident. On suppose désormais que $y$ n'est pas la fonction nulle. Il s'agit alors simplement d'une récurrence en utilisant que la dérivée de $y$ ne peut pas s'annuler en un point où $y$ s'annule, et en évaluant le wronskien en $nT$ pour obtenir l'annulation de $z$ en $nT$ donc de $y$ en $(n+1)T$. On "propage" en quelque sorte les annulations grâce au wronskien.
\end{proof}


\section{Étude d'équations différentielles d'ordre $n$}
\begin{ex}{$\Delta$ Résolution d'une EDL scalaire homogène à coefficients constant d'ordre $n$}
	On s'intéresse à l'équation différentielle scalaire homogène à coefficients constants $(E)$ : $$x^{(n)}+\sum_{k=0}^{n-1}a_kx^{(k)}=0$$
	\begin{enumerate}
		\item Montrer que, pour $\lambda\in\C$, la fonction $f_\lambda :t\mapsto e^{\lambda t}$ si, et seulement si, $\lambda$ est racine du polynôme : $$P=X^n+\sum_{k=0}^{n-1}a_kX^k$$
		\item En déduire que si le polynôme $P$ est scindé à racines simples $\lambda_1,\ldots,\lambda_n$, alors l'ensemble $E_0$ des solutions de $(E)$ est $$E_0=\text{Vect}(f_{\lambda_1},\ldots,f_{\lambda_n})$$
		\item $\star$ Décrire $E_0$ dans le cas général.
	\end{enumerate}
\end{ex}
\begin{proof}
	La dernière question est difficile.
	\begin{enumerate}
		\item Faire la vérification.
		\item La première question donne $n$ solutions indépendantes (vecteurs propres associés à des valeurs propres distinctes pour la dérivation).
		\item Utiliser le lemme des noyaux. Ce sont les $t\mapsto t^ke^{\lambda t}$ où $0\leqslant k\leqslant m_P(\lambda)-1$ avec $m_P(\lambda)$ la multiplicité de $\lambda$ en tant que racine de $P$.
	\end{enumerate}
\end{proof}
\begin{ex}{$\star$ Pas trop d'annulations dans un voisinage d'un point} On se donne une EDL d'ordre $n$ : $$x^{(n)}+\sum_{k=0}^{n-1}a_k(t)x^{(k)}=0$$ où les $a_k$ sont continues sur un intervalle $I$ d'intérieur non vide (sinon le résultat n'a pas d'intérêt), ainsi que $t_0\in I$. Alors, il existe un voisinage $J$ de $t_0$ dans $I$ tel que toute solution non nulle de l'équation s'annule au plus $n-1$ fois sur $J$.
\end{ex}
\begin{proof}
	Exercice difficile. On raisonne par l'absurde, et on prend, pour tout $p\in\N$, une fonction $\varphi_p$ solution non nulle qui s'annule au moins $n$ fois sur $$J_p=\left[t_0-\frac{1}{2^p},t_0+\frac{1}{2^p}\right]$$ Quitte à renormaliser $\varphi_p$, on peut la supposer de norme $1$ (nous verrons pour quelle norme après). En effet, cet intervalle est un voisinage de $t_0$ APCR. Alors, avec le théorème de Rolle, pour tout $k\in\llbracket0,n-1\rrbracket$, $\varphi_p^{(k)}$ s'annule en $x_p^{(k)}$. On extrait de $(\varphi_p)$ une sous-suite qui converge vers $\varphi$ par compacité de la sphère unité en dimension finie. Alors, on choisit maintenant d'avoir travaillé avec la norme $$N(f)=\lVert f\rVert_\infty+\ldots+\lVert f^{(n-1)}\rVert_\infty$$ On a alors : $$\lvert \underbrace{\varphi_p^{(k)}(t_0)}_{=0}-\varphi^{(k)}(t_0)\rvert \leqslant \lVert \varphi_p^{(k)}-\varphi^{(k)}\rVert_\infty\leqslant N(\varphi_p-\varphi)\xrightarrow[p\to+\infty]{}0$$ donc, par continuité et unicité de la limite $\varphi^{(k)}(t_0)=0$. Par unicité de la solution au problème de Cauchy avec de telles conditions initiales, $\varphi=0$, ce qui est absurde car elle est de norme $1$ donc non nulle.
\end{proof}
\begin{ex}{$\star$ Solutions stables par dérivation $\implies$ coefficients constants}
	Soient $a_0,\ldots,a_{n-1}$ continue. On suppose que l'ensemble $E_0$ des solutions de $$x^{(n)}+\sum_{k=0}^{n-1}a_k(t)x^{(k)}=0$$ est stable par dérivation. Alors les $a_k$ sont constantes.
\end{ex}
\begin{proof}
	Exercice difficile. On note $D$ l'endomorphisme induit par la dérivation sur $E_0$. $E_0$ est de dimension $n$, on note $$\chi_D=X^n+\sum_{k=0}^{n-1}b_kX^k$$ le polynôme caractéristique de $D$. D'après le théorème de Cayley-Hamilton, tout élément de $E_0$ vérifie donc aussi l'équation différentielle : $$x^{(n)}+\sum_{k=0}^{n-1}b_kx^{(k)}=0$$ En considérant les solutions dont les conditions initiales sont $$\forall k\in\llbracket0,n-1\rrbracket,\ x^{(k)}(t_0)=\delta_{k,l}$$ on on obtient que $$\forall k\in\llbracket0,n-1\rrbracket,\ \forall t\in\R,\ a_k(t)=b_k$$
\end{proof}
\begin{rmq}
	Réciproquement, avec la même méthode, tout sous-espace vectoriel de dimension finie $n$ de $\mathcal{C}^{\infty}(I,\K)$ stable par dérivation est l'ensemble des solutions d'une équation différentielle linéaire d'ordre $n$ à coefficients constants.
\end{rmq}
\begin{ex}{$\Delta$ CNS de périodicité d'une solution d'une "équation périodique"}
	On considère l'équation différentielle d'ordre $n$ suivante : $$y^{(n)}+\sum_{k=0}^{n-1}a_k(t)y^{(k)}=b(t)$$ où les $a_k$ et $b$ sont des fonctions continues toutes périodiques de même période $T$. Soit $y$ une solution de cette équation. Montrer que $y$ est $T$-périodique si, et seulement si, $$\forall k\in\llbracket0,n-1\rrbracket,\ y^{(k)}(T)=y^{(k)}(0)$$
\end{ex}
\begin{proof}
	Considérer $z:t\mapsto y(t+T)$. $z$ est aussi solution de l'équation puisque les $a_k$ et $b$ sont $T$-périodiques. Ensuite, $y$ est périodique si, et seulement si, $y=z$, ce qui équivaut, puisqu'elles sont solutions d'une même équation différentielle, à l'égalité des dérivées $k$-èmes en $0$, et correspond bien à la CNS recherchée.
\end{proof}
\section{TD Châteaux}
\begin{ex}{Lemme de Gronwall}
	\begin{enumerate}
		\item Lemme de Gronwall : soient $u$ et $g$ deux fonctions continues positives sur $[0,+\infty[$ et $A$ un réel positif tels que $$\forall x\in\R_+,\ u(x)\leqslant A+\int_{0}^{x}u(t)g(t)\ dt$$ Montrer alors que $$\forall x\in\R_+,\ u(x)\leqslant A\exp\left(\int_{0}^{x}g(t)\ dt\right)$$ Indication : considérer $x\mapsto\displaystyle\int_{0}^{x}u(t)g(t)\ dt$. On notera que $A=0$ implique $u$ nulle.
		\item Généralisation : soient $u$, $f$ et $g$ des fonctions continues positives sur $[0,+\infty[$ telles que $$\forall x\in\R_+,\ u(x)\leqslant f(x)+\int_{0}^{x}u(t)g(t)\ dt$$ Montrer alors que $$\forall x\in\R_+,\ u(x)\leqslant f(x)+\int_{0}^{x}f(t)g(t)\exp\left(\int_{t}^{x}g(u)\ du\right)\ dt$$
	\end{enumerate}
\end{ex}
\begin{ex}{Applications du lemme de Gronwall}
	\begin{enumerate}
		\item Soit $q:\R_+\to\R$ intégrable. Montrer que toute solution sur $\R_+$ de $y''+(1+q(x))y=0$ est bornée.
		\item Soit $q:\R_+\to\R$ de classe $\mathcal{C}^1$, croissante et à valeurs strictement positives. Montrer que toute solution sur $\R_+$ de $y''+q(x)y=0$ est bornée.
		\item On suppose la fonction $t\mapsto tq(t)$ intégrable sur $[a,+\infty[$ avec $a>0$. Soit $f$ une solution de $y''+qy=0$ sur $[a,+\infty[$.
		\begin{enumerate}[label=\alph*.]
			\item Montrer que $\displaystyle x\mapsto\frac{f(x)}{x}$ est bornée. Indication : appliquer Taylor avec reste intégral entre $a$ et $x$.
			\item Montrer que $f'(x)$ admet une limite finie $\lambda$ en $+\infty$. En déduire que $\displaystyle\frac{f(x)}{x}$ admet une limite finie en $+\infty$.
			\item La fonction $x\mapsto f(x)-\lambda x$ admet-elle une limite finie en $+\infty$ ? 
		\end{enumerate}
		\item On suppose $q$ intégrable sur $\R_+$. Démontrer que les solutions de $y''+qy=0$ sur $\R_+$ ne sont pas toutes bornées. Indication : on pourra utiliser le wronskien.
		\item Soit $q$ de classe $\mathcal{C}^1$ sur $\R_+$ telle que $q'$ soit intégrable sur $\R_+$ et $q(x)$ tende vers $0$ en $+\infty$.
		\begin{enumerate}[label=\alph*.]
			\item Montrer que toute solution sur $\R_+$ de $y''+(1+q(x))y=0$ est bornée.
			\item Soit $r$ une fonction intégrable sur $\R_+$. Montrer que toute solution sur $\R_+$ de $y''+(1+q(x)+r(x))y=0$.
		\end{enumerate}
		\item Soit $p$ une fonction continue sur $\R_+$ et $f$ une solution sur $\R_+$ de $y''+y'-p(x)y=0$.
		\begin{enumerate}[label=\alph*.]
			\item Montrer qu'il existe deux réels $a$ et $b$ teks que $$\forall x\in\R_+,\ f(x)=a+be^{-x}+\int_{0}^{x}p(t)f(t)\ dt-\int_{0}^{x}e^{t-x}p(t)f(t)\ dt$$
			\item On suppose que $p$ est intégrable sur $\R_+$. Montrer que $f$ est bornée.
		\end{enumerate}
		\item Soit $f\in\mathcal{C}^1(\R^2,\R)$ $M$-lipschitzienne par rapport à la deuxième variable. On pose $g=f+\phi$ où $\phi\in\mathcal{C}^1(\R^2,\R)$ vérifie $\lvert\phi\rvert\leqslant\varepsilon$. Soient $x$ et $y$ des solutions sur $\R_+$ des problèmes de Cauchy : $$\begin{cases}
			x'=f(t,x) \\
			x(0)=x_0
		\end{cases} \ \ \text{et} \ \ \begin{cases}
			y'=g(t,y)\\
			y(0)=y_0
		\end{cases}$$ Montrer en utilisant la généralisation du lemme de Gronwall que $$\forall t\in\R_+,\ \lvert x(t)-y(t)\rvert\leqslant\lvert x_0-y_0\rvert e^{Mt}+\frac{\varepsilon}{M}\left(e^{Mt}-1\right)$$
	\end{enumerate}
\end{ex}
\begin{ex}{Zéros des solutions d'une équation particulière}
	Soient $q$, $q_1$ et $q_2$ des fonctions continues de l'intervalle $I$ dans $\R$, et $f$ une solution non nulle de $y''+qy=0$, qu'on notera $(E(q))$.
	\begin{enumerate}
		\item Montrer que l'ensemble des zéros de $f$ dans $I$ est formé de points isolés.
		\item On suppose $q_1\leqslant q_2$. Soient $f_1$ et $f_2$ des solutions respectives de $(E(q_1))$ et $(E(q_2))$. Soient $\alpha<\beta$ deux zéros consécutifs de $f_1$ (cela a du sens d'après la première question). Montrer qu'il existe $\gamma\in[\alpha,\beta]$ tel que $f_2(\gamma)=0$, et que si $q_1-q_2$ n'est pas la fonction nulle sur $[\alpha,\beta]$, alors $\gamma\in]\alpha,\beta[$.
		\item Soient $f_1$ et $f_2$ deux solutions indépendantes de $(E(q))$ sur $I$, $\alpha<\beta$ deux zéros consécutifs de $f_1$. Montrer qu'il existe un unique zéro de $f_2$ dans $]\alpha,\beta[$.
		\item Soit $m>0$ tel que $\forall x\in I,\ q(x)\geqslant m$. Montrer que tout segment de longueur $\displaystyle\frac{\pi}{\sqrt{m}}$ inclus dans $I$ contient au moins un zéro de $f$.
		\item Soit $M>0$ tel que $\forall x\in I,\ q(x)\leqslant M$. Montrer que la distance entre deux zéros de $f$ est supérieure ou égale à $\displaystyle\frac{\pi}{\sqrt{M}}$.
		\item On suppose $I=\R_+$ et $\forall x\in\R_+,\ q(x)>0$. Peut-on affirmer que $f$ possède une infinité de zéros sur $\R_+$ ?
		\item On suppose $I=\R_+$ et $q(x)=e^x$. Montrer que l'ensemble des zéros de $f$ peut être ordonné en une suite croissante $(x_n)$ strictement croissante qui tend vers l'infini et telle que $x_{n+1}-x_n$ tende vers $0$. Déterminer ensuite un équivalent de $x_n$. Montrer que $f$ est bornée.
	\end{enumerate}
\end{ex}

\end{document}